\documentclass[UTF8,a4paper,11pt,fontset=windows]{ctexart}
\usepackage[left=19mm,right=19mm,top=17mm,bottom=18mm]{geometry}
\usepackage{amsmath,amssymb,array,longtable,booktabs,needspace,fancyhdr}
\usepackage[hidelinks]{hyperref}
\setlength{\parindent}{0pt}
\setlength{\parskip}{4pt}
\linespread{1.13}
\pagestyle{fancy}\fancyhf{}\renewcommand{\headrulewidth}{0pt}
\fancyfoot[C]{\small 第\thepage 页}
\setlength{\headheight}{14pt}
\newcommand{\R}{\mathbb R}
\newcommand{\src}[1]{\par{\scriptsize 教材来源：#1。}\par}
\newcommand{\q}[2]{\par\Needspace{4\baselineskip}\textbf{#1.}\ #2}
\newcommand{\subq}[2]{\par\Needspace{2\baselineskip}（#1）#2}
\newcommand{\opts}[4]{\par\begin{tabular*}{\linewidth}{@{}p{.475\linewidth}@{\extracolsep{\fill}}p{.475\linewidth}@{}}A. #1 & B. #2 \\ C. #3 & D. #4\end{tabular*}\par}
\newcommand{\optsfour}[4]{\par\begin{tabular*}{\linewidth}{@{}p{.235\linewidth}@{\extracolsep{\fill}}p{.235\linewidth}p{.235\linewidth}p{.235\linewidth}@{}}A. #1 & B. #2 & C. #3 & D. #4\end{tabular*}\par}
\newcommand{\parttitle}[1]{\par\Needspace{4\baselineskip}{\bfseries #1}\par}
\newcommand{\work}[1]{\par\vspace*{#1}}
\newcommand{\pts}[1]{\unskip\nobreak\hfill\mbox{（#1分）}\par}
\newcommand{\ans}[2]{\par\Needspace{3\baselineskip}\textbf{第#1题}\quad #2\par}
\newcommand{\titleblock}[2]{\begin{center}{\LARGE\bfseries #1}\par\vspace{3mm}{\normalsize #2}\end{center}}
\renewcommand{\arraystretch}{1.35}

\begin{document}
\titleblock{高一数学期中教材改编练习卷}{湘教版必修第一册\qquad 考试时间120分钟\qquad 满分150分}
班级\underline{\hspace{2.6cm}}\quad 姓名\underline{\hspace{2.6cm}}\quad 学号\underline{\hspace{2.6cm}}

\parttitle{一、单项选择题（本大题共8小题，每小题5分，共40分）}
每小题给出的四个选项中，只有一项符合题目要求。
\q{1}{已知集合$A=\{1,2,4\}$，$B=\{2,3,4,6\}$，则$A\cap B=$（\quad）。}
\optsfour{$\{1,2,3,4,6\}$}{$\{2,4\}$}{$\{1,6\}$}{$\{3\}$}
\src{p.13，习题1.1第8题}
\q{2}{设$x\in\R$，则“$x=2$”是“$x^2=4$”的（\quad）。}
\opts{充分而不必要条件}{必要而不充分条件}{充要条件}{既不充分又不必要条件}
\src{p.18，例5（1）}
\q{3}{命题$p$：$\forall x\in\R$，$x^2-4x+4\ge0$。则$p$的否定是（\quad）。}
\opts{$\forall x\in\R$，$x^2-4x+4<0$}{$\exists x\in\R$，$x^2-4x+4\le0$}{$\exists x\in\R$，$x^2-4x+4<0$}{$\forall x\in\R$，$x^2-4x+4>0$}
\src{p.23，页首练习（1）}
\q{4}{已知$x>2$，则$x+\dfrac4{x-2}$的最小值是（\quad）。}
\optsfour{$4$}{$6$}{$8$}{$2\sqrt2$}
\src{p.43，习题2.1第7（1）题}
\q{5}{下列各组函数中，表示同一个函数的是（\quad）。}
\opts{$y=x-1$与$y=\dfrac{x^2-1}{x+1}$}{$y=x^2$与$y=(x+1)^2$}{$y=\sqrt{x^2}$与$y=|x|$}{$y=x$与$y=\sqrt{x^2}$}
\src{p.77，习题3.1第5题}
\q{6}{函数$f(x)=\dfrac{\sqrt{x+2}}{x^2-x-6}$的定义域是（\quad）。}
\opts{$[-2,+\infty)$}{$(-2,3)\cup(3,+\infty)$}{$[-2,3)\cup(3,+\infty)$}{$(-\infty,3)\cup(3,+\infty)$}
\src{p.77，习题3.1第4（1）题}
\q{7}{已知$f(x)$是定义在$\R\setminus\{0\}$上的奇函数，当$x>0$时，$f(x)=x^2+\dfrac1x$，则$f(-2)=$（\quad）。}
\optsfour{$\dfrac92$}{$-\dfrac72$}{$\dfrac72$}{$-\dfrac92$}
\src{p.86，习题3.2第6题}
\q{8}{若函数$f(x)=x^2+2(a-1)x+2$在区间$(-\infty,2]$上单调递减，则实数$a$的取值范围为（\quad）。}
\optsfour{$(-\infty,-1]$}{$(-\infty,-1)$}{$[-1,+\infty)$}{$[1,+\infty)$}
\src{p.82，例3}

\newpage
\parttitle{二、多项选择题（本大题共3小题，每小题6分，共18分）}
每小题有多项符合题目要求。全部选对得6分，部分选对得3分，有选错的得0分。
\q{9}{已知$a<b<0$，则下列结论正确的是（\quad）。}
\opts{$-a>-b$}{$\dfrac1a<\dfrac1b$}{$\sqrt{a^2}>\sqrt{b^2}$}{$\dfrac ba>1$}
\src{p.43，习题2.1第3题}
\q{10}{已知$f(x)=x^2-4x+1$，$x\in[0,3]$，则下列结论正确的是（\quad）。}
\opts{$f(x)$的最小值为$-3$}{$f(x)$的最大值为$1$}{$f(x)$在$[0,3]$上单调递减}{$f(x)$的值域为$[-3,0]$}
\src{p.82，练习第2（3）题；p.90，复习题三第7题}
\q{11}{已知函数$f:\R\to\R$满足$f(x+y)=f(x)+f(y)$（$x,y\in\R$），且$f(1)=2$，则（\quad）。}
\opts{$f(0)=0$}{$f(x)$是奇函数}{$f(x)$是偶函数}{$f(3)=6$}
\src{p.91，复习题三第12题}

\parttitle{三、填空题（本大题共3小题，每小题5分，共15分）}
\q{12}{若对任意实数$x$，$x^2-2x+m\ge0$恒成立，则实数$m$的取值范围是\underline{\hspace{3cm}}。}
\src{p.55，例5}
\q{13}{已知偶函数$f(x)$的定义域为$\R$，在$[0,+\infty)$上单调递减，则$f(-3)$、$f(2)$、$f(-1)$由小到大的排列为\underline{\hspace{4cm}}。}
\src{p.86，习题3.2第11题}
\q{14}{函数$f(x)=|x-1|-|x+1|$的值域是\underline{\hspace{3cm}}。}
\src{p.76，页首练习第1（2）题}

\parttitle{四、解答题（本大题共5小题，共77分）}
解答应写出文字说明、证明过程或演算步骤。
\q{15}{（15分）设全集$U=\R$，集合$A=\{x\mid x\le2\}$，$B=\{x\mid x>a\}$，其中$a\in\R$。}
\subq{1}{当$a=-1$时，求$A\cap B$、$A\cup B$及$\complement_U A$。}
\subq{2}{若$A\cup B=\R$，求实数$a$的取值范围。}
\src{p.13，习题1.1第9、16题}
\work{4.2cm}

\newpage
\q{16}{（15分）解下列关于$x$的不等式，其中$m\in\R$。}
\subq{1}{$2x^2-x-3\le0$；}
\subq{2}{$x^2-(m+2)x+2m<0$。}
\src{p.59，习题2.3第4、8题}
\work{6.2cm}
\q{17}{（17分）已知函数$g(x)$的定义域为$[-3,3]$，设$f(x)=g(x)+g(-x)$，且当$x\in[0,3]$时，$f(x)=x^2-2x$。}
\subq{1}{判断$f(x)$的奇偶性，并证明你的结论；}
\subq{2}{求$f(x)$在$[-3,0)$上的解析式；}
\subq{3}{求$f(x)$在$[-3,3]$上的最大值、最小值，以及取得相应最值时$x$的所有取值。}
\src{p.84，例5；p.82，练习第2（3）题}
\work{7.2cm}

\newpage
\q{18}{（17分）完成下列问题。}
\subq{1}{计算$27^{\frac23}+16^{-\frac34}$的值；}
\subq{2}{已知幂函数$f(x)=x^\alpha$的图像经过点$\left(2,\dfrac14\right)$，求$f(x)$的解析式及定义域；}
\subq{3}{判断（2）中函数的奇偶性，写出其单调区间，并用定义证明它在$(-\infty,0)$上的单调性。}
\src{p.104，习题4.1第1、4、8题}
\work{6.2cm}
\q{19}{（13分）已知函数$f(x)=a^{|2x-1|}$（$a>0$且$a\ne1$），且$f(1)=\dfrac14$。}
\subq{1}{求$a$的值；}
\subq{2}{写出$f(x)$的单调区间及对应的增减性，并求其值域；}
\subq{3}{解不等式$f(x)\ge\dfrac1{16}$。}
\src{p.114，习题4.2第11题}
\work{6.8cm}
\end{document}
